\documentclass[12pt,letterpaper]{article}
\usepackage[margin=1in,headheight=15pt]{geometry}
\usepackage{newtxtext,newtxmath}
\usepackage{setspace}
\usepackage[section]{placeins}
\usepackage[font=normalsize,labelfont=bf]{caption}
\usepackage[T1]{fontenc}
\usepackage[utf8]{inputenc}
\usepackage{amsmath,graphicx,booktabs,array}
\usepackage{microtype}
\usepackage{xurl}
\usepackage[hidelinks]{hyperref}
\urlstyle{tt}
\newcommand{\tightlist}{\setlength{\itemsep}{0pt}\setlength{\parskip}{0pt}}
\setlength{\emergencystretch}{1.5em}
\interdisplaylinepenalty=2500
\raggedbottom
\graphicspath{{figures/}}
\hypersetup{pdftitle={Comparing Hyperparameter Search Methods on Tabular Data},pdfauthor={Optima Prime}}
\renewcommand{\thesection}{\Roman{section}}
\renewcommand{\thesubsection}{\Alph{subsection}}
\makeatletter
\renewcommand{\@seccntformat}[1]{\csname the#1\endcsname.\quad}
\makeatother
\renewcommand{\thetable}{\Roman{table}}
\renewenvironment{abstract}{\section*{Abstract}}{\par}
\pagestyle{plain}
\setlength{\parindent}{0.25in}
\setlength{\parskip}{0pt}
\doublespacing
\begin{document}
\thispagestyle{empty}
\begin{center}
\singlespacing
\vspace*{1.1in}
{\LARGE\bfseries Budgeted Hyperparameter Optimization on Tabular Data: Predictive Performance, Computational Cost, and Robustness\par}
\vspace{0.65in}
{\Large MAP 6114 Machine Learning\par}
\vspace{0.45in}
{\large\bfseries Team: Optima Prime\par}
\vspace{0.3in}
{\large Danay Fernandez Alfonso\par}
\vspace{0.1in}
{\large Dilek Ozdemirci\par}
\vspace{0.1in}
{\large Ecil Teodoro\par}
\end{center}
\clearpage
\setcounter{page}{1}
\title{\bfseries\fontsize{18}{22}\selectfont Budgeted Hyperparameter Optimization on Tabular Data: Predictive Performance, Computational Cost, and Robustness}
\author{Danay Fernandez Alfonso, Dilek Ozdemirci, and Ecil Teodoro\\
\normalsize Team Optima Prime, MAP 6114 Machine Learning}
\date{}
\begingroup
\singlespacing
\maketitle
\endgroup

\section{Introduction}\label{introduction}

Tuning hyperparameters often takes much longer than training a single model. A regularization strength, kernel setting, or tree setting that works well on one dataset can work badly on another. Our team wanted to see when a simple search is good enough and when an adaptive method is worth the extra effort.

To look at this, we compared four search methods on six real datasets and four synthetic ones. Together they include regression and classification, small and large sample sizes, few and many features, and balanced and imbalanced classes. We also changed the training data in two ways, by adding noise and by using fewer samples, to see how the methods held up. Within each comparison, every optimizer got the same data split, the same cross-validation folds, and the same seed.

Our main questions were \textbf{which method gives better test predictions}, \textbf{what happens when the budget goes from 25 to 100}, and \textbf{how long each method takes}. We also looked at how much the results moved between seeds and whether our conclusions changed under other metrics. The results we have answer some of these questions, but not all of them, since some planned comparisons are still missing.

The computing time forced us to change the plan. We went from ten seeds to five, which cut the planned optimization runs from 3,060 to 1,530. Even so, the main five-seed group used 135.72 hours of recorded search time without being complete. Run one after another, that is 5.65 days. If we include the older code versions, the total goes up to 164.70 hours, or 6.86 days. Everything ran on a 16-core Apple M4 Max with 128 GB of RAM, and we kept five-fold cross-validation and both budgets.

There is one more difference from our proposal. The proposal described giving each method the same amount of time. What we actually did was limit the number of full evaluations (or a rough equivalent for halving) and record the elapsed time separately. This matters, since two methods can use about the same number of evaluations and still take very different amounts of time.

\section{Related Work}\label{related-work}

\subsection{Random and Bayesian search}\label{random-and-bayesian-search}

Bergstra and Bengio \cite{bergstra12} show why random search can be a better choice than a grid. If only a few hyperparameters really affect performance, a grid wastes a lot of trials changing the ones that don't. This was our reason for including random search as a baseline, especially for LightGBM, which has a bigger search space.

Bergstra et al.~\cite{bergstra11} describe sequential methods that use past results to pick the next configuration to try, including the approach behind TPE. For our Bayesian optimization runs we used Optuna's TPE sampler. So our results are about this specific implementation and not about every method that gets called Bayesian optimization.

\subsection{Allocating resources during a search}\label{allocating-resources-during-a-search}

Jamieson and Talwalkar \cite{jamieson16} treat hyperparameter search as a problem of finding a good candidate while also deciding how much computing each candidate should get. This is the idea behind successive halving: try many candidates cheaply, then give more resources to the ones that look good. Whether it actually saves time depends in part on how well the early scores predict the final ones.

\subsection{Tuning, evaluation, and model choice}\label{tuning-evaluation-and-model-choice}

Bischl et al.~\cite{bischl21} cover search algorithms, but also the practical side of tuning: evaluation, preprocessing, parameter ranges, and parallel runs. The point we took most directly from them was about preprocessing. If the median used to fill missing values or the mean used for standardizing is computed on the whole training set, the validation fold has already ``seen'' part of itself, and the CV score looks better than it should. So in our code, every fold computes these values from its own training part only. We also kept the test set completely out of parameter selection and saved the search settings for every run.

Demšar \cite{demsar06} discusses how to compare several learning algorithms statistically over many datasets. We based our planned Friedman and post-hoc tests on this paper, and it is also why we did not treat different seeds as independent datasets. We have not reported formal test results yet, because our coverage and convergence checks are not met.

Ke et al.~\cite{ke17} present LightGBM, a fast gradient-boosted tree method. Its large parameter space makes it a useful contrast to the linear models and SVM. Optuna \cite{akiba19} provides our TPE search, and Scikit-learn \cite{pedregosa11} provides the other estimators, preprocessing, folds, and scoring functions. Grinsztajn, Oyallon, and Varoquaux \cite{grinsztajn22} study model performance on tabular data while taking tuning and model assumptions into account. Their work supports comparing several model families, but it also reminds us that the model type and the size of the search space change together.

\section{Methodology}\label{methodology}

\subsection{Datasets and experiment design}\label{datasets-and-experiment-design}

We planned 19 dataset-condition combinations and 51 dataset-condition-model combinations, which we call cells. Linear models and LightGBM have 19 cells each. SVM has 13, because we left it out of the Base and Noisy conditions for the larger datasets. Every cell has four optimizers, budgets of 25 and 100, and five seeds. In total the design needs 1,530 optimization runs and gives 2,040 budget summaries. There are fewer runs than summaries because random search and TPE use one search history for both budgets.

\begin{table}[t]
\centering
\caption{Datasets and training conditions. Row and feature counts are source values before cleaning.}
\label{tab:datasets}
\fontsize{10}{12}\selectfont
\setlength{\tabcolsep}{3.3pt}
\begin{tabular}{lrrrr}
\toprule
Dataset (OpenML ID) & Task & Rows & Features & Conditions \\
\midrule
White Wine (44971) & Reg. & 4,898 & 11 & Base, Noisy \\
Superconductivity (44964) & Reg. & 21,263 & 81 & Base, Small, Noisy \\
Communities/Crime (46286) & Reg. & 1,994 & 122 & Base, Noisy \\
Bank Marketing (1461) & Class. & 45,211 & 16 & Base, Small, Noisy \\
Bioresponse (4134) & Class. & 3,751 & 1,776 & Base, Noisy \\
Credit Default (42477) & Class. & 30,000 & 23 & Base, Small, Noisy \\
Four synthetic generators & Both & 2,000 each & 20 & Base \\
\bottomrule
\end{tabular}
\end{table}

Table I gives the row and feature counts of the original datasets, before our cleaning and encoding. For Bank Marketing, we dropped call duration, since it is only known after the call is over. For Communities and Crime, we dropped the identifiers and a fixed list of problem attributes, then converted everything left to numbers. For Credit Default, we dropped the identifier and treated some integer-coded columns as categories. Because of these steps, the number of inputs a model actually sees can be different from the source feature count. The table is not meant as a new audit of every dataset property.

For each seed, we first held out 20\% of the data as a test set. Classification splits were stratified so the class proportions stayed the same. The Base condition uses all the remaining training data. The Small condition cuts the training set down to 2,000 observations, sampled with stratification by class or by quantiles of the regression target. We only used Small for Superconductivity, Bank Marketing, and Credit Default.

In the Noisy condition, we only changed the training labels. For regression, we added Gaussian noise with standard deviation \(0.5\,s_{y,\text{train}}\), where \(s_{y,\text{train}}\) is the standard deviation of the training target. For classification, we flipped 10\% of the training labels. For a given dataset and seed, Base, Small, and Noisy all share the same clean test set, so comparing them shows the effect of changing only the training data.

Each synthetic task has 2,000 observations and 20 features. The linear regression generator has ten informative features and a target-noise standard deviation of 10. The Friedman generator has five informative features and a noise standard deviation of 1. The classification generators have five informative and five redundant features, a label-reassignment rate of 0.01, and either about equal classes or a nominal 90:10 split. We generated a new dataset for each seed. The synthetic tasks only use Base, and the generator's built-in noise is added before the split, so it is in both the training and test data.

\subsection{Models and preprocessing}\label{models-and-preprocessing}

For regression we used Elastic Net, and for classification we used elastic-net logistic regression with the SAGA solver. We call both of these the linear-model family. The SVM uses an RBF kernel, with probability estimates turned on for classification. LightGBM gives us gradient-boosted tree models for both regression and classification \cite{ke17}. Because SVM is slow to train, we only ran it on the smaller datasets and the Small conditions.

For numeric features, we filled missing values with the median and then standardized. For categorical features, we filled missing values with the most common category and used one-hot encoding that can handle categories it has not seen before. These steps were fit separately inside each cross-validation training fold, so the validation fold never affected the preprocessing.

LightGBM used the same preprocessing. We thought about using its built-in categorical handling, but did not end up implementing it. The models also differ in more ways than just the number of parameters we tune. That means we can compare how the optimizers behave across model families, but we can't explain every difference by the size of the search space alone.

\begin{table}[t]
\centering
\caption{Hyperparameter ranges used by the searches. For halving, LightGBM tree count is set by the resource stage.}
\label{tab:spaces}
\fontsize{10}{12}\selectfont
\setlength{\tabcolsep}{3.3pt}
\begin{tabular}{p{0.23\textwidth}p{0.70\textwidth}}
\toprule
Model / group & Parameters and bounds \\
\midrule
Linear / regression & alpha: 0.0001-100 (log); l1\_ratio: 0.05-1 \\
Linear / classification & C: 0.0001-100 (log); l1\_ratio: 0-1; class\_weight: none/balanced \\
SVM / both & C: 0.001-1000 (log); gamma: 0.00001-1 (log) \\
SVM / task-specific & epsilon: 0.001-1 (log, regression); class\_weight: none/balanced (classification) \\
LightGBM / structure & num\_leaves: 8-128; max\_depth: -1, 3, 5, 8, 12; min\_child\_samples: 5-100 \\
LightGBM / sampling & subsample: 0.5-1; colsample\_bytree: 0.5-1 \\
LightGBM / boosting & learning\_rate: 0.01-0.3 (log); n\_estimators: 50-500 \\
LightGBM / regularization & reg\_alpha and reg\_lambda: 0.000001-10 (log); min\_split\_gain: 0-1 \\
\bottomrule
\end{tabular}
\end{table}

The linear models have two or three tunable parameters, SVM has three, and LightGBM has ten in the full setup. Table II lists the ranges, and values marked logarithmic are sampled on a log scale. Elastic Net for regression can run up to 50,000 iterations, and the classification version up to 10,000. We saved convergence warnings instead of automatically throwing out those candidates. Because of this, some of the selected models still have warnings. We report how many, and we include a separate table that only uses comparisons where neither the selected fit nor the refit had a warning.

\subsection{Search methods and budgets}\label{search-methods-and-budgets}

For grid search, our code builds the largest rectangular grid that fits inside the budget. This has a side effect that matters later in the results: if a parameter ends up with only one value on the grid, the code just takes the first category or the lower end of the range, and that isn't always the model's default. Random search is simpler and just draws from the ranges in Table II. TPE is the only one of the three that looks back at the validation scores so far before choosing its next candidate \cite{bergstra11}, \cite{akiba19}. In all four methods the trials run one after another, so the only parallel work happens inside a trial, where the five CV folds are fit at the same time.

Our successive-halving code starts with randomly sampled candidates and uses three resource levels, \(r \in \{1/9, 1/3, 1\}\). After each stage, roughly a third of the candidates move on. For linear models and SVM, the resource is the size of the training fold. The subsets are nested, and the validation folds do not change. For LightGBM, the resource is the number of boosting rounds, which reaches the maximum of 500 in the last stage. At each stage we fit the candidates from scratch instead of continuing from the earlier fit. Only candidates that were evaluated at full resource can be picked for the final refit.

We count how much resource halving uses as

\[B_{\text{used}} = \sum_{j=1}^{J} r_j.\]

Here, one full evaluation means testing one configuration on all five CV folds, not fitting one model once. The saved halving schedules use 24 and 99 full-evaluation equivalents for the nominal budgets of 25 and 100. The LightGBM grids test 24 and 96 configurations. Every LightGBM grid configuration selected in this snapshot has 50 trees. So even with the same nominal budget, the methods do not do exactly the same amount of fitting or cover the parameters in the same way.

Random search and TPE each run for 100 trials. To get their budget-25 result, we take the best of the first 25 trials, and we refit the chosen configuration separately for each budget. Grid and halving are run separately at each budget, because their search designs change with the budget. The two summaries that come from one shared history are related to each other, and adding up their times would count part of the same search twice.

\subsection{Model selection, scores, and hardware}\label{model-selection-scores-and-hardware}

We used the same five cross-validation folds for every optimizer, so all four searches were judged on exactly the same splits. Regression was scored with RMSE and classification with ROC-AUC. If you open our output files, you'll notice RMSE is saved as a negative number during the search. Scikit-learn always maximizes its scores, so it flips the sign of error metrics; we flip it back in the summary tables. After the search picked a winner, we retrained it on all of the training data and scored it once on the clean test set.

We also saved mean absolute error (MAE) for regression, and F1, average precision (AP), and log loss for classification. F1 combines precision and recall using each model's default decision threshold. AP summarizes how precision and recall trade off as the threshold changes. It is not the trapezoid-rule area under a precision-recall curve. Log loss looks at the predicted probabilities and penalizes confident wrong answers.

For comparisons where lower should always be better, we define the test loss as

\[\ell = \begin{cases} \text{RMSE}, & \text{regression},\\ -\text{ROC-AUC}, & \text{classification}. \end{cases}\]

The gain from the bigger budget is \(g = \ell_{25} - \ell_{100}\). A positive gain means the budget-100 model did better on the test set. We define the generalization gap as test RMSE minus validation RMSE, or validation AUC minus test AUC, so a positive gap means the test result was worse. In the Noisy condition, validation uses the changed training labels and testing uses the clean labels, so the gap there is not purely a measure of overfitting.

All experiments ran on a 16-core Apple M4 Max with 128 GB of RAM. The main analysis group used four parallel CV workers with one native model thread each. We did not check whether all 16 cores were actually busy. The recorded optimization time includes the search and CV work, but not downloading data or the final refits. Search overhead, fitting cost, and other programs running on the machine can all affect elapsed time. We did not measure CPU usage or energy, and multiplying elapsed hours by the number of cores would not give a real CPU-hour total. Software versions are listed in Appendix A.

\subsection{Matching runs and comparing methods}\label{matching-runs-and-comparing-methods}

We only include a comparison when all four optimizers have results for the same dataset, condition, model, budget, and seed. Within that group, we rank their test losses from 1 to 4, and ties get the average rank. So that datasets with more conditions don't count more, we first average the seed ranks within each condition, then average the conditions within each dataset. The four synthetic generators are closely related, so we treat them as one synthetic-family block. The final model-budget averages are taken over these dataset blocks. We never average raw RMSE across datasets, since their targets are on different scales.

Our planned statistics are Friedman tests across blocks, a Holm correction across the six model-budget tests, and Nemenyi comparisons when the corrected omnibus test is significant \cite{demsar06}. Before running the tests on a panel, we require all planned cells and seeds, passing convergence checks for the selected trials and refits, and at least five blocks. Right now, no panel meets all of these. Because of this, we do not report any p-values or claim statistical significance. Even if the study were complete, there would only be seven benchmark blocks, so the tests would have limited power and the large-sample approximations would have to be used carefully.

We also compare test performance and time using Pareto efficiency. In a matched group, a method is Pareto efficient if no other method is at least as good on both test loss and time while being strictly better on one of them. This shows which combinations of speed and accuracy were useful in our results. It does not let anyone pick a method for new data by looking at that data's test set.

Some plots show differences relative to the absolute value of a reference score. A relative change in AUC is not the same as a change in AUC percentage points, and it also does not mean the same thing as a relative change in RMSE. We use these plots to describe the results, not as a shared effect-size scale.

\section{Timeline and Teamwork Plan}\label{timeline-and-teamwork-plan}

We held our Week 6 meeting on Zoom on September 26, 2026. As of the October 30 snapshot, the repository has dataset loaders, preprocessing inside the CV folds, model factories, and all four search methods. It also has atomic checkpoints so saved results can't get corrupted, reuse of runs that already finished, convergence diagnostics, and a fully executed analysis notebook. The work was split among the three team members as described below.

\begin{table}[t]
\centering
\caption{Completed milestones and the proposed schedule for finishing the project.}
\label{tab:timeline}
\fontsize{10}{12}\selectfont
\setlength{\tabcolsep}{3.3pt}
\begin{tabular}{p{0.20\columnwidth}p{0.71\columnwidth}}
\toprule
Period & Milestone / next action \\
\midrule
Week 6 & Scope and team coordination meeting (documented). \\
Weeks 7-9 & Runner, saved experiments, and executed analysis implemented. \\
Week 10 & Check report claims and submit the midterm PDF with cover letter. \\
Weeks 11-12 & Resolve feasible coverage/convergence issues and review grid/resource semantics. \\
Weeks 13-14 & Refresh a frozen snapshot; review results, limitations, and references. \\
Week 15 & Finalize the paper and team presentation. \\
\bottomrule
\end{tabular}
\end{table}

Table X shows the project milestones and the rest of the schedule. Ecil Teodoro did the coding. Danay Fernandez Alfonso chose the datasets and analyzed the results. Dilek Ozdemirci analyzed the results and wrote the documentation.

\begin{thebibliography}{11}
\interlinepenalty=10000
\bibitem{bergstra12} J. Bergstra and Y. Bengio, "Random Search for Hyper-Parameter Optimization," Journal of Machine Learning Research, vol. 13, pp. 281-305, 2012. \href{https://jmlr.org/papers/v13/bergstra12a.html}{Online}.
\bibitem{bergstra11} J. Bergstra, R. Bardenet, Y. Bengio, and B. Kegl, "Algorithms for Hyper-Parameter Optimization," Advances in Neural Information Processing Systems, vol. 24, 2011. \href{https://proceedings.neurips.cc/paper/2011/hash/86e8f7ab32cfd12577bc2619bc635690-Abstract.html}{Online}.
\bibitem{jamieson16} K. Jamieson and A. Talwalkar, "Non-stochastic Best Arm Identification and Hyperparameter Optimization," Proceedings of AISTATS, PMLR, vol. 51, pp. 240-248, 2016. \href{https://proceedings.mlr.press/v51/jamieson16.html}{Online}.
\bibitem{bischl21} B. Bischl et al., "Hyperparameter Optimization: Foundations, Algorithms, Best Practices and Open Challenges," arXiv:2107.05847v3, 2021. \href{https://arxiv.org/abs/2107.05847}{Online}.
\bibitem{demsar06} J. Demsar, "Statistical Comparisons of Classifiers over Multiple Data Sets," Journal of Machine Learning Research, vol. 7, pp. 1-30, 2006. \href{https://jmlr.org/papers/v7/demsar06a.html}{Online}.
\bibitem{ke17} G. Ke et al., "LightGBM: A Highly Efficient Gradient Boosting Decision Tree," Advances in Neural Information Processing Systems, vol. 30, 2017. \href{https://proceedings.neurips.cc/paper/2017/hash/6449f44a102fde848669bdd9eb6b76fa-Abstract.html}{Online}.
\bibitem{akiba19} T. Akiba, S. Sano, T. Yanase, T. Ohta, and M. Koyama, "Optuna: A Next-generation Hyperparameter Optimization Framework," Proceedings of KDD, 2019; arXiv:1907.10902. \href{https://arxiv.org/abs/1907.10902}{Online}.
\bibitem{pedregosa11} F. Pedregosa et al., "Scikit-learn: Machine Learning in Python," Journal of Machine Learning Research, vol. 12, pp. 2825-2830, 2011. \href{https://jmlr.org/papers/v12/pedregosa11a.html}{Online}.
\bibitem{grinsztajn22} L. Grinsztajn, E. Oyallon, and G. Varoquaux, "Why do tree-based models still outperform deep learning on tabular data?" Advances in Neural Information Processing Systems, vol. 35, pp. 507-520, 2022. \href{https://arxiv.org/abs/2207.08815}{Online}.
\end{thebibliography}

\end{document}